


% Please add the following required packages to your document preamble:
% \usepackage{booktabs}
% \usepackage{multirow}
% \usepackage[normalem]{ulem}
% \useunder{\uline}{\ul}{}
\begin{table*}[]
\centering \small
\addtolength{\tabcolsep}{-1.0mm}
\begin{tabular}{@{}llccccccccccc@{}}
\toprule
\multirow{2}{*}{\textbf{Model}} & \multirow{2}{*}{\textbf{Feature}} & \multirow{2}{*}{\textbf{Length}} & \multicolumn{2}{c}{\textbf{Rouge-1 (F1)}} & \multicolumn{2}{c}{\textbf{Rouge-2 (F1)}} & \multicolumn{2}{c}{\textbf{Rouge-L (F1)}} & \multicolumn{2}{c}{\textbf{MoverScore}} & \multicolumn{2}{c}{\textbf{BERTScore}} \\ \cmidrule(l){4-13}
 &  &  & \textbf{Score} & \textbf{Norm} & \textbf{Score} & \textbf{Norm} & \textbf{Score} & \textbf{Norm} & \textbf{Score} & \textbf{Norm} & \textbf{Score} & \textbf{Norm} \\ \midrule
\multirow{6}{*}{\textbf{Reuse}} & \textbf{M} & 33.2 & .291 & 1.346 & .139 & 1.790 & .239 & 1.401 & .535 & 1.023 & .628 & 1.064 \\
 & \textbf{P} & 238.3 & .171 & 1.042 & .089 & 1.006 & .134 & 1.030 & .503 & 1.004 & .567 & 1.008 \\
 & \textbf{W{[}0, 1{]}} & 50.3 & .281 & 1.216 & .132 & 1.509 & .224 & 1.273 & .529 & 1.016 & .620 & 1.048 \\
 & \textbf{W{[}0, 2{]}} & 68.0 & .259 & 1.129 & .123 & 1.341 & .205 & 1.186 & .524 & 1.013 & .611 & 1.034 \\
 & \textbf{W{[}1, 1{]}} & 67.8 & .266 & 1.156 & .124 & 1.346 & .204 & 1.183 & .524 & 1.012 & .613 & 1.037 \\
 & \textbf{W{[}2, 2{]}} & 98.7 & .235 & 1.082 & .112 & 1.179 & .180 & 1.105 & .517 & 1.007 & .600 & 1.020 \\ \midrule
\multirow{6}{*}{\textbf{Pegasus}} & \textbf{M} & 12.2 & .321 & 1.898 & .153 & 2.907 & .283 & 1.971 & .553 & 1.065 & .654 & 1.158 \\
 & \textbf{M+O} & 12.8 & .331 & 1.909 & .161 & 2.945 & .292 & 1.993 & .556 & 1.071 & .661 & 1.166 \\
 & \textbf{P} & 14.0 & {\ul .374} & {\ul 2.067} & {\ul .205} & {\ul 3.507} & {\ul .334} & {\ul 2.201} & {\ul .570} & {\ul 1.095} & {\ul .682} & {\ul 1.196} \\
 & \textbf{P+O} & 14.0 & \textbf{.381} & \textbf{2.106} & \textbf{.212} & \textbf{3.635} & \textbf{.340} & \textbf{2.242} & \textbf{.571} & \textbf{1.097} & \textbf{.685} & \textbf{1.202} \\
 & \textbf{P+O+B} & 38.3 & .321 & 1.452 & .154 & 1.916 & .265 & 1.537 & .546 & 1.044 & .639 & 1.082 \\
 & \textbf{O} & 12.1 & .133 & 0.789 & .026 & 0.495 & .119 & 0.828 & .518 & 0.998 & .561 & 0.993 \\ \midrule
\textbf{TrOCR} & \multirow{2}{*}{\textbf{Figure}} & 10.0 & .220 & 1.464 & .073 & 1.653 & .195 & 1.502 & .534 & 1.033 & .610 & 1.096 \\
\textbf{BEiT+GPT2} &  & 15.8 & .164 & 0.864 & .042 & 0.666 & .144 & 0.917 & .529 & 1.013 & .592 & 1.031 \\
\bottomrule
\end{tabular}
\addtolength{\tabcolsep}{+1.0mm}
\caption{Task Performance with the \textbf{best} and {\ul second-best} results highlighted. Pegasus$_{P+O}$, the text-summarization model with all available information (Paragraph+OCR), performed the best in all four metrics. Pegasus$_{P+O+B}$, the model trained with better captions, however, got lower scores.}
\label{tab:result-summarization}
\vspace{-3mm}
\end{table*}


\begin{comment}
% Please add the following required packages to your document preamble:
% \usepackage{booktabs}
% \usepackage{multirow}
\begin{table*}[]
\centering \small

\addtolength{\tabcolsep}{-1.2mm}

\begin{tabular}{@{}l@{\kern3pt}l@{\kern5pt}rrrrrrrrrrrrr@{}}
\toprule
\multirow{2}{*}{\textbf{Model}} & \multirow{2}{*}{\textbf{Feature}} & \multicolumn{1}{c}{\multirow{2}{*}{\textbf{Length}}} & \multicolumn{3}{c}{\textbf{ROUGE-1 (F1)}} & \multicolumn{3}{c}{\textbf{ROUGE-2 (F1)}} & \multicolumn{3}{c}{\textbf{ROUGE-L (F1)}} & \multicolumn{3}{c}{\textbf{MoverScore}} \\ \cmidrule(l){4-15} 
 &  & \multicolumn{1}{c}{} & \multicolumn{1}{c}{\textbf{Score}} & \multicolumn{1}{c}{\textbf{Rand}} & \multicolumn{1}{c}{\textbf{Norm}} & \multicolumn{1}{c}{\textbf{Score}} & \multicolumn{1}{c}{\textbf{Rand}} & \multicolumn{1}{c}{\textbf{Norm}} & \multicolumn{1}{c}{\textbf{Score}} & \multicolumn{1}{c}{\textbf{Rand}} & \multicolumn{1}{c}{\textbf{Norm}} & \multicolumn{1}{c}{\textbf{Score}} & \multicolumn{1}{c}{\textbf{Rand}} & \multicolumn{1}{c}{\textbf{Norm}} \\ \midrule
\multirow{6}{*}{\textbf{Reuse}} & \textbf{M} & 33.2 & .291 & .216 & 1.346 & .139 & .077 & 1.790 & .239 & .171 & 1.401 & .535 & .524 & 1.023 \\
 & \textbf{P} & 238.3 & .171 & .164 & 1.042 & .089 & .088 & 1.006 & .134 & .130 & 1.030 & .503 & .501 & 1.004 \\
 & \textbf{W{[}0, 1{]}} & 50.3 & .281 & .231 & 1.216 & .132 & .087 & 1.509 & .224 & .176 & 1.273 & .529 & .520 & 1.016 \\
 & \textbf{W{[}0, 2{]}} & 68.0 & .259 & .230 & 1.129 & .123 & .092 & 1.341 & .205 & .173 & 1.186 & .524 & .518 & 1.013 \\
 & \textbf{W{[}1, 1{]}} & 67.8 & .266 & .230 & 1.156 & .124 & .092 & 1.346 & .204 & .173 & 1.183 & .524 & .518 & 1.012 \\
 & \textbf{W{[}2, 2{]}} & 98.7 & .235 & .217 & 1.082 & .112 & .095 & 1.179 & .180 & .162 & 1.105 & .517 & .513 & 1.007 \\ \midrule
\multirow{6}{*}{\textbf{Pegasus}} & \textbf{M} & 12.2 & .321 & .169 & 1.898 & .153 & .053 & 2.907 & .283 & .144 & 1.971 & .553 & .519 & 1.065 \\
\textbf{} & \textbf{M+O} & 12.8 & .331 & .173 & 1.909 & .161 & .055 & 2.945 & .292 & .147 & 1.993 & .556 & .520 & 1.071 \\
\textbf{} & \textbf{P} & 14.0 & .374 & .181 & 2.067 & .205 & .058 & 3.507 & .334 & .152 & 2.201 & .570 & .521 & 1.095 \\
\textbf{} & \textbf{P+O} & 14.0 & \textbf{.381} & .181 & \textbf{2.106} & \textbf{.212} & .058 & \textbf{3.635} & \textbf{.340} & .152 & \textbf{2.242} & \textbf{.571} & .521 & \textbf{1.097} \\
\textbf{} & \textbf{P+O+B} & 38.3 & .321 & .221 & 1.452 & .154 & .080 & 1.916 & .265 & .172 & 1.537 & .546 & .523 & 1.044 \\
\textbf{} & \textbf{O} & 12.1 & .133 & .168 & 0.789 & .026 & .052 & 0.495 & .119 & .143 & 0.828 & .518 & .519 & 0.998 \\ \midrule
\textbf{TrOCR} & \multirow{2}{*}{\textbf{Figure}} & 10.0 & .220 & .150 & 1.464 & .073 & .044 & 1.653 & .195 & .130 & 1.502 & .534 & .517 & 1.033 \\
\textbf{BEiT+GPT2} &  & 15.8 & .164 & .190 & 0.864 & .042 & .062 & 0.666 & .144 & .158 & 0.917 & .529 & .522 & 1.013 \\ \bottomrule
\end{tabular}

\addtolength{\tabcolsep}{+1.2mm}

\caption{Task Performance. Pegasus$_{P+O}$, the text-summarization model with all available information (Paragraph+OCR), performed the best in all four metrics. Pegasus$_{P+O+B}$, the model trained with better captions, however, got lower scores.}
\label{tab:result-summarization}
\vspace{-2mm}

\end{table*}

\end{comment}


\begin{comment}


\begin{table*}[t]
    \centering
    %\small
    % \addtolength{\tabcolsep}{-0.02cm}
    
    \caption{Task Performance. The text-summarization model (Pegasus) with all available information (Paragraph+OCR) performed the best in all three metrics. The model with the same setting that was trained with better captions, however, got lower scores. \cy{Should we add ROUGE-L?}}
    \vspace{-2mm}
    \label{tab:result-summarization}
    \begin{tabular}{llrrrrrrrrrrr}
    \toprule 
    % \hline
    \multirow{2}{*}{\textbf{Model}} & \multirow{2}{*}{\textbf{Feature}} & \multicolumn{1}{c}{\multirow{2}{*}{\textbf{Length}}} & \multicolumn{3}{c}{\textbf{ROUGE-1 (F1)} $\uparrow$ }  & \multicolumn{3}{c}{\textbf{ROUGE-2 (F1)} $\uparrow$ } & \multicolumn{3}{c}{\textbf{MoverScore} $\uparrow$} \\
    % \cline{4-12}
    \cmidrule{4-12}
     & &  & \multicolumn{1}{c}{\textbf{Score}} & \multicolumn{1}{c}{\textbf{Rand}} & \multicolumn{1}{c}{\textbf{Norm}} & \multicolumn{1}{c}{\textbf{Score}} & \multicolumn{1}{c}{\textbf{Rand}} & \multicolumn{1}{c}{\textbf{Norm}} & \multicolumn{1}{c}{\textbf{Score}} & \multicolumn{1}{c}{\textbf{Rand}} & \multicolumn{1}{c}{\textbf{Norm}} \\ 
     \midrule
    %  \hline
\multirow{6}{*}{\textbf{Reuse}} & \textbf{Mention}      & 33.2  & .2912 & .2164 & 1.346 & .1387 & .0775 & 1.790 & .5354 & .5236 & 1.023 \\ 
& \textbf{Paragraph}                        & 238.3 & .1706 & .1638 & 1.042 & .0887 & .0882 & 1.006 & .5030 & .5009 & 1.004 \\ 
& \textbf{Window [0, 1]}                    & 50.3  & .2808 & .2309 & 1.216 & .1318 & .0874 & 1.509 & .5290 & .5205 & 1.016 \\ 
& \textbf{Window [0, 2]}                    & 68.0  & .2595 & .2298 & 1.129 & .1231 & .0918 & 1.341 & .5242 & .5175 & 1.013 \\ 
& \textbf{Window [1, 1]}                    & 67.8  & .2657 & .2298 & 1.156 & .1236 & .0918 & 1.346 & .5237 & .5176 & 1.012 \\ 
& \textbf{Window [2, 2]}                    & 98.7  & .2349 & .2171 & 1.082 & .1116 & .0947 & 1.179 & .5170 & .5132 & 1.007 \\ 

\midrule

\multirow{6}{*}{\textbf{Pegasus}} & \textbf{Mention}             & 12.2  & .3207 & .1690 & 1.898 & .1533 & .0528 & 2.907 & .5528 & .5191 & 1.065 \\ 
& \textbf{Mention+OCR}         & 12.8  & .3310 & .1733 & 1.909 & .1613 & .0548 & 2.945 & .5563 & .5196 & 1.071 \\ 
& \textbf{Paragraph}           & 14.0  & .3745 & .1812 & 2.067 & .2047 & .0584 & 3.507 & .5700 & .5206 & 1.095 \\ 
& \textbf{Paragraph+OCR}       & 14.0  & \textbf{.3811} & .1810 & \textbf{2.106} & \textbf{.2118} & .0583 & \textbf{3.635} & \textbf{.5712} & .5205 & \textbf{1.097} \\ 
& \textbf{Paragraph+OCR-Better} & 38.3  & .3205 & .2207 & 1.452 & .1541 & .0804 & 1.916 & .5455 & .5227 & 1.044 \\ 
& \textbf{OCR}                 & 12.1  & .1328 & .1682 & 0.789 & .0259 & .0524 & 0.495 & .5178 & .5190 & 0.998 \\ 

\midrule

\textbf{BEiT+GPT2} & \multirow{2}{*}{\textbf{Figure}} & 15.8 & .1642 & .1901 & 0.864 & .0415 & .0624 & 0.666 & .5287 & .5217 & 1.013 \\
\textbf{TrOCR} &  & 10.0 & .2196 & .1500 & 1.464 & .0735 & .0444 & 1.653 & .5341 & .5168 & 1.033 \\ 

\bottomrule
    \end{tabular}
    %\vspace{-2mm}
    % \addtolength{\tabcolsep}{0.02cm}
    % \vspace{-.5pc}
    
\end{table*}
\end{comment}

%In this section, we first describe the experimental setups and then show our automatic evaluation and human evaluation results.\kenneth{Opening needs a rewrite}

% \paragraph{Dataset.}
% %\gromit{I feel sudden when jumping from related work to data. Feels like we assume people know our method needs data and annotations. I add a sentence below to make a transition.}

\kenneth{Add a few sentences to say (1) the remaining of the paper was built using SciCap, including motivating analysis and testing of the proposed methods and in-depth analysis. (2) We additionally add mention/paragraph and resplit the data (suggesting we did a lot more than just using it.)}

\kenneth{I agree with Gromit that it reads a bit sudden. We can add a sentence to explain that we want to introduce dataset here as it was used in Motivating Anlaysis too.}
\cy{done}

%The figure-captioning framework requires a model to be trained from a large corpus with figures and texts. 
%Before introducing our method, we will describe the data used in the experiments.

Before diving into our experiments and analyses, we first describe the dataset upon which our study is grounded.
Our results are based on a scientific figure caption dataset, \oldDataset, and several pre-processing steps to fit it into our workflow.

% \footnote{We will release the code and data upon publication. Please see 200 data samples in the supplementary material (\url{https://anonymous.4open.science/r/Generating-Figure-Captions-as-a-Text-\\Summarization-Task-15B4}).}
%\kenneth{The footnote should be AFTER the period or ,}
%\paragraph{\oldDataset}
%
%\paragraph{Extracting Mentions and Paragraphs for Figures in \oldDataset.}
% \paragraph{Extracting Mentions and Paragraphs for Figures in \oldDataset.}
%
\oldDataset is a dataset that contains over 416,000 line charts and captions extracted from more than 290,000 arXiv papers~\cite{hsu-etal-2021-scicap-generating}. 
It was one of the first large-scale figure-captioning datasets based on real-world \sfigs. 
However, it does not contain the paragraphs that mention the figure.
To address this, we downloaded all the PDF files of the original arXiv papers used in \oldDataset and
extracted all the Mentions and Paragraphs as outlined in \Cref{sec:extract-mention}.
% We then followed the steps described in Section~\ref{sec:extract-mention} to extract all the Mentions and Paragraphs.
% It is worth noting that we manually verified 399 figures (the set used in \Cref{sec:analysis}) and found that 81.2\% (324/399) were published at academic conferences, and 51.9\% (207/399) were at ACL Anthology, IEEE, or ACM, suggesting that the data is representative.\kenneth{``It is worth noting...that the data is representative.'' can be moved to Appendix.}
Detailed information on preprocessing, including the dataset resplit and OCR extraction, are described in \Cref{sec:appendix-data-preprocessing}.

\begin{comment}

In this paper, ``\textbf{Mention}'' refers to the sentence in an article that explicitly refers to the target figure, and 
``\textbf{Paragraph}'' refers to all the text under the same \texttt{<p>} tag (paragraph tag) as the Mention in the PDF parsing result.
As sentences near a Mention might be more relevant, we also extracted $n$ preceding sentences and $m$ following sentences to form \textbf{Windows~[n, m]}.
For example, ``Windows [1, 2]'' refers to a text snippet of five sentences, \ie, one preceding sentence, one Mention sentence, and two following sentences.

% \vspace{-1pc}

\end{comment}















%------------------ dead kitten ----------------

\begin{comment}




Five models were trained using five different inputs, 
including Mention, Paragraph, OCR, Mention+OCR, and Paragraph+OCR. 
Paragraph+OCR covered all the relevant information in this paper and thus should provide the best summary. 


%Pre-training with Extracted Gap-sentences for 

the Pegasus text-summarization model~\cite{zhang2020pegasus}, fine-tuned on our dataset. 
Five models were trained using five different inputs, 
including Mention, Paragraph, OCR, Mention+OCR, and Paragraph+OCR. 
Paragraph+OCR covered all the relevant information in this paper and thus should provide the best summary. 
One special variation of the model, Paragraph+OCR-Better, was intended to be trained with better-quality captions. 
We selected sample captions longer than 30 tokens to be \textbf{better-quality} samples (27,224 samples remaining) 
because caption length is positively correlated to caption quality (see \Cref{sec:quality-annotation}).\kenneth{Any justification for picking 30?} 
Prior work also argued that increasing caption length could increase readers' comprehension~\cite{hartley2003single,gelman2002let}.
This model was then trained with Paragraph+OCR inputs on only the better-quality subset.

\kenneth{The cutting on 30 tokens need to be better justified.}





% \paragraph{OCR}
\subsection{OCR}
% \cy{TODO: ED please add the tool you use.}
We used EasyOCR~\cite{easyocr}
% \footnote{EasyOCR: \url{https://github.com/JaidedAI/EasyOCR}} 
to extract texts from all the figures.
EasyOCR provided OCR texts along with their bounding boxes.
%To represent the OCR texts into a single sentence, 
When feeding OCR texts into models, 
we joined the OCR texts with the order of the coordinates from left to right and then top to bottom.

\begin{itemize} %[leftmargin=*,itemsep=0mm]

    \item 
    A ``\textbf{Mention}'' refers to a sentence in an article that explicitly mentions the target figure, such as ``As shown in Figure 6 ...'' Since there might be multiple Mentions throughout the paper, without further specification, we referred to the first Mention.
    
%\vspace{-.5pc}

    \item 
    
    A ``\textbf{Paragraph}'' refers to all the text under the same \texttt{<p>} tag as the Mention in the PDF-parsing result.

    
%\vspace{-1pc}

    \item 
    As sentences near a Mention may also be relevant, we extracted $n$ preceding sentences and $m$ following sentences to form \textbf{Window[n, m]}. For example, ``Window[1, 2]'' refers to a text snippet of four sentences, \ie, one preceding sentence, one Mention sentence, and two following sentences.

\end{itemize}

\bigskip

used Grobid~\cite{grobid}, a publicly available tool that parses PDF files into structured XML documents, 
to extract the plain text of all the paragraphs (including the \texttt{<p>} tags) in each paper. 
Next, we used BlingFire~\cite{blingfire} to segment each sentence. 
We created regular expressions to identify the sentences that mention specific figures. 
For example, the sentence ``As shown in Figure 6, ...'' would be identified and linked to Figure 6. 
To this end, we used the following terminology:

%Paper : 53771/6721/6722
%Figure : 106391/13926/13226
%CL figure: 3559 / 419 / 438 (train/val/test)

%
%\gromit{What is \oldDataset? needs a brief intro for new readers.}
%\ed{We also notice that the figures are sometimes merely captioned and not mentioned in the paper. We then further filter out the figures that are not mentioned in the paper because there are no text references to do summary. This step is do before resplit}


% \kenneth{We only find mentions for 80\% of images. I moved this to Limitations, which gives us some more space.}
\oldDataset was originally created for vision-to-language tasks and we need new train/val/test splits for our work.
As figures do not overlap with each other, different figures in the same paper could be assigned to different data splits in \oldDataset.
However, paper's texts overlap with each other more easily and can be 
%As \oldDataset assigned figures to each split randomly,
%In \oldDataset, a paper's figures could be assigned to different data splits.
%This would be reasonable for vision-to-language tasks, , 
%but can 
problematic for text summarizing tasks. %, as texts can overlap with each other. 
We resplit \oldDataset to make sure no paper has figures from different data splits; we also excluded the figures whose mentions without any identified mentions.
As a result, 
the train/val/test sets used in the paper have
86,825/10,833/10,763 figures, 
sourced from 48,603/6,055/6,053 papers.






Extracted mentions, ...
Get better result with mentions.
Next step: use not only the first sentence, but ones after, as they are more likely to add additional info about the insight, key result, etc.



\kenneth{TODO Ed: Fill in the data details.}
\cy{We only used samples whose Mention exists. This is also different from \oldDataset.}

\paragraph{Data Source.}
\ed{We download the PDF from Arxiv dataset. We then use Grobid, a public ML library to parse PDF into structured XML encoded documents. We only keep the plain text in <p> tag that represents each paragraph in the paper.} 




% Figure Caption with Pretrained OCR models (baseline)
% Also, had some examples where OCR would improve (small captions, or bad captions, cumulative...)

\paragraph{Text Pre-processing.}

\ed{We retrieve the figure index from the figure caption (Fig. 1, Figure 2, etc.) and then use regular expression to discover all the paragraphs in the paper that mention corresponding figure index. BlingFire is also used to separate the sentences in the paragraph. The sentence mention the figure index that we called it \textbf{Mention}. To better understand the relationship between caption and mention, we also remove the figure index in the caption when we do experiment and analysis.}
 
\cy{TODO: ED please describe (1) how we identify identify Mentions, (2) remove figure index}


\paragraph{Extracting Mentions and Paragraphs.}
% Probably give a definition for both ``Mention'', ``Window[n, m]'', and``Paragraph''.
\cy{Should we show an example here?}
The \textbf{Mention} refers to the sentence in the paper that mentions the target figure (e.g., the sentence before (Fig. X)). 
The \textbf{Paragraph} is all the text under the same \texttt{<p>} tag (paragraph tag) 
with the Mention in the PDF parsing result.
As sentences closer to Mention might be more relevant, 
we also extract $n$ preceding sentences and $m$ following sentences to form \textbf{Windows[n, m]}.
For example, Windows [1, 2] will result in a text snippet of five sentences (one preceding sentence, one Mention sentence, and two following sentences).
To extract the Mention we use regular expression. 
\cy{TODO: ED please explain how we get the mention.}


\end{comment}






%\subsection{Baselines}

\paragraph{A Simple Baseline: Using Extracted Mentions as Captions.}
%\kenneth{MAYBE not put it here...}
%As suggested by prior literature that images usually illustrate objects in the surrounding documents~\cite{6205758} and 
Motivated by our information overlap study (\Cref{sec:motivating-analysis}),
we created the \textbf{Reuse} baselines. 
These baselines simply repurpose portions of the input text as the prediction.
%we designed the \textbf{Reuse} baselines that simply output parts of the input text (\eg, texts surrounding Mention) as the prediction. 
% For example, Reuse$_M$ (Reuse with Mention) outputs the first mention as the caption;
% Reuse$_{W[0, 1]}$ (Reuse with Window[0, 1]) outputs the first mention and the following 1 sentence as the caption.
% \kenneth{CY TODO: Check if it makes sense.}
% and Reuse$_{P}$ (Reuse with Paragraph) outputs the whole paragraph as the summary.
%Similar to prior findings that relevant events are more likely to occur closely~\cite{huang-huang-2021-semantic}, 
%we expected Reuse$_M$ to perform the best among the Reuse baselines and its performance would drop as the window size increases.

% \kenneth{We need to cite some more papers to justify this baseline. (1) Semantic Frame Forcast paper and (2) maybe this \cite{6205758}}
% \cy{Done}

% \paragraph{Baseline: Vision-to-Language Generation.}
\paragraph{Vision-to-Language Baselines.}
The vision-to-language generation treated this task as an image-captioning task 
that took the scientific figure image as input and generated a text to describe it.
%For this baseline, 
We compared two vision-to-language models as baselines.
First, we built a sequence-to-sequence model by combining BEiT~\cite{bao2022beit} and GPT-2~\cite{radford2019language}.
%The second model we chose was the TrOCR~\cite{li2021trocr} model, a transformer-based sequence-to-sequence model pre-trained for the OCR task.
%Compared to image encoders trained on real photos (images containing visual objects) 
%such as ViT~\cite{dosovitskiy2020vit} and BEiT~\cite{bao2022beit}, 
%OCR models trained with images of printed and handwritten documents are closer to the scientific paper domain, 
%where figures are mostly lines and texts.
We also selected the TrOCR~\cite{li2021trocr} model, a transformer-based sequence-to-sequence model pre-trained for OCR tasks.
Compared to image encoders like ViT~\cite{dosovitskiy2020vit} and BEiT~\cite{bao2022beit}, which were trained on photos, OCR models trained on printed and handwritten documents align more closely with the scientific paper domain.
All figures from \oldDataset (106,391 training samples) were used for training since no mentions were required.
%Note that 
%All available figures from \oldDataset (106,391 training samples) were used for training, as no mentions were needed.
%\kenneth{Is is before or after resplit?}\cy{After the resplit.}
% HuggingFace's implementation of the TrOCR model\footnote{\texttt{microsoft/trocr-large-printed} was used.} was used for fine-tuning.
% The model was trained for 100 epochs, and the one with the highest ROUGE-2 score was kept.

% \paragraph{Baseline: [X-Axis] and [Y-Axis]}

% \kenneth{TODO ED}

\begin{figure*}[t]
    \centering
    \includegraphics[width=0.48\linewidth]{figures/length-performance-rouge2-fmeasure.png}
    \includegraphics[width=0.48\linewidth]{figures/length-performance-rouge2-fmeasure-normalized.png}
    \vspace{-3mm}
    \caption{The relationship between average text length and ROUGE-2 score (left: original ROUGE-2; right: normalized ROUGE-2). The random baseline in the left figure shows that text lengths and scores are not independent. For example, when the predicted text is shorter than 50 tokens, predicting longer texts generally results in a higher ROUGE-2 score. The normalized scores indicate the proposed system's performance gain over the random baseline of the same length. Pegasus$_{P+O+B}$ and Reuse$_{M}$ get closer to TrOCR after normalization, suggesting the need for normalization for accurate interpretation of results.} 
    \label{fig:length-performance}
    \vspace{-3mm}
\end{figure*}

\paragraph{Experimental Setup.}
\label{sec:training-setup}
% \kenneth{TODO TY: If possible, replace the footnotes with citations. Footnotes eat space.}

A total of 14 methods were included for comparison: 
six reuse baselines with six input variations
% (Mention, Paragraph, Window[0, 1], Window[0, 2], Window[1, 1], and Window[2, 2]); 
(M, P, W[0, 1], W[0, 2], W[1, 1], and W[2, 2]);
five text summarization models with five inputs
% (Mention, Mention+OCR, Paragraph, Paragraph+OCR, and OCR); 
(M, M+O, P, P+O, and O);
one text summarization model using P+O with controlled data quality; 
and two vision-to-language models (BEiT+GPT-2 and TrOCR). 
Note that we use subscripts of M, P, W, O, B to denote the input features: 
Mention, Paragraph, Window, OCR, and Better data quality, respectively.
% For example, Reuse$_{W[0, 1]}$ refers to the Reuse baseline outputting Window[0, 1] and Pegasus$_{P+O+B}$ refers to the text summarization model using better quality of Paragraph+OCR as the input.
% \kenneth{We could probably remove FOR EXAMPLE to save space.}
The model training details and decoding configuration are provided in \Cref{sec:appendix-model-training-details}.
% Next, we describe the model training details for both the text summarization models and the vision-to-language models.
% Reuse$_{M}$, Reuse$_{P}$, Reuse$_{W[0, 1]}$
% Pegasus$_{M}$, Pegasus$_{M+O}$, Pegasus$_{M+O+B}$

% \paragraph{Model Training Detail.}\label{sec:appendix-model-training}
% \cy{Need update}

% \kenneth{TODO CY: Move the following details to Appendix.}

%We then applied both automatic evaluation and human evaluation to compare model performance.





\subsection{Automatic Evaluation Results}
\paragraph{Conventional Automatic Evaluation.}
We used F1 of ROUGE-1, ROUGE-2, ROUGE-L~\cite{lin-2004-rouge,nallapati-etal-2016-abstractive}, 
MoverScore~\cite{zhao-etal-2019-moverscore}, and BERTScore for automatic evaluation. 
% ROUGE-1 F1 and ROUGE-2 F1 evaluated the F1 score of the unigram and bigram 
% token overlapping predicted texts and ground-truth captions. 
% ROUGE-L F1 computed the F1 score based on the longest common subsequence instead.
When computing ROUGE scores using rouge-score~\cite{rouge-scorer},
%\sukim{if the definition is clear and consistent, we don't have to cite this}
%\footnote{We used rouge-score: \url{https://github.com/google-research/google-research/tree/master/rouge}}
we turned all text into lower case and stem words.
As both MoverScore and BERTScore are based on the semantic similarity,
we obtained contextual embeddings from SciBERT~\cite{beltagy-etal-2019-scibert}.
%\footnote{We used \texttt{allenai/scibert\_scivocab\_uncased}.}

% MoverScore evaluated semantic similarity between predicted texts and ground-truth captions by 
% (\textit{i}) extracting contextual embeddings for words using BERT-based models and 
% (\textit{ii}) computing the Word Mover's Distance (WMD)~\cite{kusner2015word}, 
% where similar words are aligned and the cost for traveling between these words is accumulated. 
% In this work, we obtained contextual embeddings from SciBERT~\cite{beltagy-etal-2019-scibert}.\footnote{We used \texttt{allenai/scibert\_scivocab\_uncased}.}
% \cy{Add BERTScore}




% \begin{figure}
%     \centering
%     \includegraphics[width=0.9\linewidth]{figures/length-performance-rouge1 - fmeasure.png}
%     \vspace{-3mm}
%     \caption{ROUGE-1 scores with respect to length. Random shows within 50 tokens; longer texts lead to a higher score.}
%     \label{fig:rouge-1-length-performance}
%     \vspace{-2mm}
% \end{figure}



\paragraph{Automatic Evaluation with Normalization Over Caption Length.}
%ROUGE F1 has a limitation in that its scores favor longer texts within a certain length range, even though they take both precision and recall into account~\cite{sun-etal-2019-compare}, resulting in a misleading comparison, 
%as models that generate longer texts get higher scores. 
ROUGE F1 tends to favor longer texts within a certain length, leading to a skewed comparison where models generating longer texts receive higher scores~\cite{sun-etal-2019-compare}.
%Therefore, it is necessary to normalize automatic evaluation metrics with respect to the length of the text. 
We followed \citet{sun-etal-2019-compare}'s approach of normalizing the scores 
with the corresponding random baseline that generates texts of the same length. 
\begin{equation}
    \small
    Score_{normalized} = \frac{Score}{Random(length)}
\end{equation}
where $length$ is the average length of the texts generated by the target system.
We estimated $Random(length)$ by applying linear interpolation on several (length, random score) pairs.
The (length, random score) pairs were obtained by randomly selecting a certain number of sentences (1, 2, ..., 10 sentences) from the input paragraph as the prediction. 
To get random scores of texts shorter than a single sentence (around 30 tokens), 
we truncated sentences to the desired length (4, 6, ..., 30 tokens). 
For each length setting, we ran 10 different random seeds and reported the average. 
The Random line in \Cref{fig:length-performance}-Left shows the behavior of ROUGE-2
favoring longer texts within 50 tokens.\footnote{Similar trends for ROUGE-1, ROUGE-L, MoverScore, and BERTScore 
% can be seen in \Cref{fig:length-performance-rouge1} to \Cref{fig:length-performance-bertscore} 
are included in \Cref{sec:additional-results}.}
% \Cref{fig:length-performance-rouge1}, \Cref{fig:length-performance-rougel}, \Cref{fig:length-performance-wms}, \Cref{fig:length-performance-bertscore} in \Cref{sec:additional-results}).
% \footnote{ROUGE-1 shows the same behavior (see \Cref{sec:additional-results}).} 
The normalized scores, as shown in \Cref{fig:length-performance}-Right, 
%intuitively show how much better the proposed model is than the random baseline.
clearly indicate the superiority of our proposed model over the random baseline.


\Cref{tab:result-summarization} shows the normalized automatic evaluation results. 
Overall, \textbf{Pegasus$_{P+O}$, the text-summarization model with all available information (Paragraph+OCR), 
achieved the best performance in all three metrics.}
Pegasus$_{P+O+B}$, the model using the same information but trained on a better subset of captions (Paragraph+OCR-Better), did not perform well. 
%We hypothesized that this was because half of the testing data actually consisted of bad captions (see \Cref{sec:quality-annotation}). 
%We verified this by examining performance changes in beams of different quality (\Cref{sec:auto-eval-quality-beam}) and applying human evaluation to the ensuing paragraphs (\Cref{sec:human-eval}). 
We hypothesized this was due to half of the test data comprising poor captions (refer to \Cref{sec:quality-annotation}). 
This was validated by examining performance shifts in different quality beams (\Cref{sec:auto-eval-quality-beam}) and conducting a human evaluation (\Cref{sec:human-eval}).
Meanwhile, Reuse$_{M}$, the Reuse baseline with Mention, outperforms other Reuse baselines.
Its performance declined as context sizes grew and shifted.
%Reuse$_{M}$, the Reuse baseline with Mention, is the strongest Reuse baseline.
%The performance dropped when the context sizes increased and moved.
%implying that the surrounding contexts of Mention contained more relevant information. 
%We also found that OCR (Pegasus$_{O}$) was ineffective even compared with the vision-to-language models (TrOCR and BEiT+GPT2). 
%This may suggest that replacing pre-processed OCR texts with representations from a vision encoder would further improve performance.
%\kenneth{I removed the Pegasus$_{O}$ vs. OCR. Did we mention this later in the paper? If no, we can remove it. }

% \cy{say more about the intuitive of the normalization. How much better you are compared to the random baseline.}
% Interpolation
% 









% \begin{figure*}[t]
%     \centering
%     \includegraphics[width=0.98\linewidth]{figures/Quality-Beam-Normalized-Rouge1.png}
%     \caption{Normalized Rouge-1 for beams of different quality. \cy{We can probably move this figure to Appendix.}}
%     \label{fig:quality-beam-rouge-1}
% \end{figure*}








\subsection{Human Evaluation Results\label{sec:human-eval}}

%\kenneth{Add the small MTurk study we mentioned in the rebuttal.}

\paragraph{Pilot MTurk Study to Select Top Models.}
\cy{Probably emphasize why we choose these four systems for comparison}
Before the main human evaluation, we ran a pilot study on Amazon Mechanical Turk (MTurk) to identify any apparently underperforming baselines for exclusion in the final study, simplifying the main human evaluations.
In this study, we asked MTurk workers to carefully read a figure and select the \textit{worst} figure caption among 
{\em (i)} TrOCR, 
{\em (ii)} Pegasus$_{P+O}$, 
{\em (iii)} Pegasus$_{P+O+B}$, and 
{\em (iv)} ground-truth caption. 
%\sukim{The notation for the models above should be consistently used, including when you first introduce them at the beginning of section 4.2.}
Ninety figures without errors were randomly sampled from our annotated set (\ie, figures from \texttt{cs.CL} arXiv papers in \Cref{sec:quality-annotation}) for the study.
For each of the figures, we recruited 20 MTurk workers to judge.\footnote{Four MTurk qualifications were used: Locale (US Only), HIT Approval Rate ($\geq$98\%), Number of Approved HITs ($\geq$3000), and the Adult Content Qualification. The payment for each task was set to 0.09 (hourly wage = \$10 dollars).}
We report the number of majority votes (when tied, we counted all captions with the highest votes as the worst) and the average number of votes in \Cref{tab:mturk-eval}.
Results indicated that TrOCR's caption won the majority vote 41 out of 90 times, with its average vote count significantly exceeding others.
Hence, we excluded TrOCR from our formal human evaluation.


\paragraph{Main Human Evaluation with Domain Experts.}
%The human evaluation was conducted to assess the quality of the generated captions. 
Three Ph.D. students with NLP backgrounds (who are not coauthors) were recruited as human judges, 
as it is hard for those without basic domain understanding to evaluate captions.
This study has been approved by the IRB office of the authors' institute.
The same 90 figures used in the pilot MTurk study were used again.
% We randomly sampled 90 figures without errors from our annotated set (\ie, figures from \texttt{cs.CL} arXiv papers in \Cref{sec:quality-annotation}) 
We asked the human judges to compare each figure's 
(\textit{i}) Pegasus$_{P+O}$, (\textit{ii}) Pegasus$_{P+O+B}$, and (\textit{iii}) ground-truth caption. 
The judges were asked to rank the captions based on how strongly they agreed with this statement:
``When I read the paper, this caption can help me understand the message that the figure tries to convey.'' 
% Figure~\ref{fig:ui-compare} in \Cref{sec:app-interface} shows the interface and instructions for human judges.
\Cref{fig:ui-compare} (see \Cref{sec:app-interface}) shows the interface the human judges used. %where the initial order of the captions was randomized to avoid bias.
%The paper's title (without linking to the paper's URL) and abstract are shown. 
%The human judges can drag the captions (each displayed with the figure) 
%from the left pane into to the right pane to rank them. 
%The initial display order of the captions is randomized on the interface. 
%To prevent human judges from biasing toward the author-written captions, we did not display the paper's PDF or link to the paper's URL. 


% Please add the following required packages to your document preamble:
% \usepackage{multirow}
\begin{table}[]
\centering \small
\begin{tabular}{@{}l@{}c@{\kern2pt}c@{\kern2pt}c@{\kern2pt}c@{\kern2pt}c@{}}
\toprule
 \multirow{2}{*}{\textbf{n = 90}} & \multirow{2}{*}{\makecell{\textbf{\#Maj.} \\ \textbf{Votes}$\downarrow$}} & \multirow{2}{*}{\makecell{\textbf{Avg.} \\ \textbf{Votes}$\downarrow$}} & \multicolumn{3}{c}{\textbf{T-Test over Avg. Votes}} \\ \cmidrule{4-6}
 &  &  & \textbf{Peg}$_{P+O}$ & \textbf{Peg}$_{P+O+B}$ & \textbf{Caption} \\ \midrule
\textbf{TrOCR} & 41 & 5.99 & \textless .001*** & .006** & .001** \\
\textbf{Peg}$_{P+O}$ & 20 & 4.54 & - & .253 & .973 \\
\textbf{Peg}$_{P+O+B}$ & 24 & 4.93 & - & - & .318 \\
\textbf{Caption} & 19 & 4.53 & - & - & - \\
\bottomrule
\end{tabular}
\caption{The result of the pilot Mturk study. When tied, all captions with the highest votes were counted as the worst for \#Majority votes. TrOCR is significantly worse than other approaches when rated by crowd workers.}
\label{tab:mturk-eval}
\vspace{-3mm}
\end{table}




\begin{comment}
& \multirow{2}{*}{\textbf{\#Majority Vote}$\downarrow$} & \multirow{2}{*}{\textbf{Avg. \#Vote}$\downarrow$} & \multicolumn{3}{c}{\textbf{T-Test over Avg. \# Vote}} \\ \cmidrule{4-6}
 &  &  & \textbf{Pegasus}$_{P+O}$ & \textbf{Pegasus}$_{P+O+B}$ & \textbf{Caption} \\ \midrule

& \textbf{\#Majority} & \textbf{Avg.} & \multicolumn{3}{c}{\textbf{T-Test over Avg. \#Vote}} \\ \cmidrule{4-6}
 & \textbf{Vote}$\downarrow$ & \textbf{\#Vote}$\downarrow$ & \textbf{Pegasus}$_{P+O}$ & \textbf{Pegasus}$_{P+O+B}$ & \textbf{Caption} \\ \midrule

\end{comment}


% Please add the following required packages to your document preamble:
% \usepackage{multirow}
\begin{table}[]
    \centering \small
    
    \begin{tabular}{@{}lccc@{}}
    \toprule
    \multirow{2}{*}{\textbf{n = 90}} & \multirow{2}{*}{\makecell{\textbf{Avg.} \\ \textbf{Ranking}$\downarrow$}} & \multicolumn{2}{c}{\textbf{T-Test on Avg. Ranking}} \\ \cmidrule{3-4}
     &  & \textbf{Peg}$_{P+O+B}$ & \textbf{Caption} \\ \midrule
    \textbf{Peg}$_{P+O}$ & 2.152 & .016* & .015* \\
    \textbf{Peg}$_{P+O+B}$ & 1.930 & - & .923 \\
    \textbf{Caption} & \textbf{1.919} & - & - \\ \bottomrule
    \end{tabular}
    
    \caption{Average ranking of the human evaluation. Pegasus$_{P+O+B}$ was rated significantly better than Pegasus$_{P+O}$ and was at the same level as the ground-truth caption.}
    \vspace{-2mm}
    \label{tab:human-eval}
    \vspace{-4mm}
\end{table}






\begin{comment}
\multirow{2}{*}{\textbf{Avg. Ranking} $\downarrow$} 

	generation-long	generation	original	size
all	        1.930	2.152	1.919	270 (90)
helpful	    1.970	2.176	1.855	165 (55)
unhelpful	1.867	2.114	2.019	105 (35)



\begin{table}[t]
    \centering
    \small
    \caption{Average ranking of human evaluation in different quality beams. Pegasus$_{P+O+B}$ was considered the most helpful when the original caption was bad.}
    \vspace{-2mm}
    \label{tab:human-eval}
    % \addtolength{\tabcolsep}{-0.14cm}
    \begin{tabular}{lrrr}
    \toprule 
    % \hline
    $\downarrow$ & \textbf{Annotated} & \textbf{Helpful} & \textbf{Unhelpful} \\ 
     \midrule
    %  \hline
    \textbf{\# Samples}              & 90 & 55 & 35 \\ 
    \midrule
    % \hline
    \textbf{Pegasus}$_{P+O}$           & 2.152 & 2.176 & 2.114 \\
    \textbf{Pegasus}$_{P+O+B}$         & 1.930 & 1.970 & \textbf{1.867} \\
    \textbf{Caption}                 & \textbf{1.919} & \textbf{1.855} & 2.019 \\
    % \hline 
    \bottomrule
    \end{tabular}
    % \addtolength{\tabcolsep}{+0.14cm}
    \vspace{-3mm}
    
\end{table}

\end{comment}

\begin{comment}

As shown in Figure~\ref{fig:ui-compare}, 
we instructed human judges to 
(\textit{i}) rank the caption as \textit{bad} if all information is obviously incorrect, and 
(\textit{ii}) ignore errors and focus on ``whether the caption can help understanding the figure'' if only a small part of the information is incorrect.\kenneth{TODO: UPDATE the REFs}

\end{comment}


\Cref{tab:human-eval} shows the results of average ranking (from 1 to 3). 
Overall, \textbf{the ground-truth caption and Pegasus$_{P+O+B}$ were ranked similarly} (1.919 vs. 1.930 with p-value = 0.923). 
\textbf{Humans also favored Pegasus$_{P+O+B}$ over Pegasus$_{P+O}$ significantly} (1.919 vs. 2.152 with p-value = 0.016).
This supports our heuristic for automatically determining caption quality based on length 
and aligns with previous findings that longer captions improve reader comprehension~\cite{hartley2003single,gelman2002let}. 
However, we found that the task of caption ranking poses a challenge, as evidenced by the lower correlations between raters, with Kendall's tau values of 0.133, 0.148, and 0.274, and Spearman's rho values of 0.128, 0.156, and 0.317.
This highlights the complexity of the task and suggests that scaling human evaluation across domains might be difficult.
Different preferences over captions, such as length, could lead to lower agreement among raters.

% However, we found that the caption ranking task is challenging as evidenced by the lower correlations between raters (Kendall's tau: 0.133, 0.148, and 0.274; and Spearman's rho: 0.128, 0.156, 0.317.)
% The values emphasize the complexity of the task and suggest that it might be difficult to scale the human evaluation to various domains.
% Another possible hypothesis is that people might have different preferences over captions, leading to low agreements.
% (short concise but correct information v.s. rich information but small number of errors).




%------------------- dead kitten -------------------

\begin{comment}

%------------ old draft below -----------

%------------ old draft below -----------




%\cy{A really rough draft. Need to rewrite.}\kenneth{Will take a look. Should probably add an UI screenshot to appendix.}
We conduct human evaluation to further examine the quality of the generated captions. Because figure captions vary greatly, it is difficult for those without basic domain understanding to evaluate them. We then recruit three Ph.D. students with NLP background to evaluate the quality.
We randomly sample 90 samples that has no errors from our annotated set (test set captions from {\tt cs.CL} papers). 
We compare Pegasus-Paragraph+OCR, Pegasus-Paragraph+OCR-Better, and ground truth caption. 
The task is designed to be a ranking task where annotators are asked to rank the captions to answer the question ``When I read the paper, this caption can help me understand the message that the figure tries to convey.''
In our pilot study, we notice that the generated captions sometimes contain factual errors.
Therefore, we instruct annotators to 
(\textit{i}) rank the caption as a bad caption if all the information is obviously incorrect;
(\textit{ii}) ignore the error and focus on ``whether the caption can help understanding the figure'' if only a small part of the information is wrong. \ed{maybe discuss more in error analysis}

\Cref{tab:human-eval} shows the results of average ranking (from one to three). 
Again, we provide the scores in beams of different quality to examine the behavior.
Overall, ground truth caption and Pegasus-Par.+OCR-BTR are ranked similarly (1.919 v.s 1.930) and both of them are significantly better than Pegasus-Par.+OCR (2.152) with p-values of 0.0145 and 0.0165, respectively. 
Considering beams of different quality, we find that Pegasus-Par.+OCR-BTR (1.867) is ranked better than ground truth captions (2.019) when they are not helpful. 

% We extracted CL papers 
% CL figure: 3559 / 419 / 438 (train/val/test)



Again, we presented the scores in beams of different quality to examine the behavior. 
Overall, the ground-truth caption and Pegasus$_{P+O+B}$ were ranked similarly (1.919 vs. 1.930 with p-value = 0.923)
and were both significantly better than Pegasus$_{P+O}$ (2.152), with p-values of 0.0145 and 0.0165, respectively. 
When considering beams of different quality, 
Pegasus$_{P+O+B}$ (1.867) was ranked better than ground-truth captions (2.019) when they were not helpful. 
On average, human judges preferred machine-generated captions 22 out of 35 times. 
%As a result, ZZ\% of the captions were considered not helpful to readers.
%We also found that the NLP researchers often preferred machine-generated captions, \ie, X out of Y times, for these cases where the original caption is unhelpful.
%Finally, we investigated the relations between the captions and the paragraphs that mention them.
The results showed an early evidence that, with some quality controls over training data, 
machines can summarize relevant paper contents to generate better captions and help readers.
However, due to the cost of recruiting experts to assess system outputs, 
more research is needed to perform a comprehensive evaluation of machine-generated captions on a larger scale.
%, especially when the human-written ones are bad.
%\kenneth{Maybe worth one more closing sentence.}
%\cy{done.}
%%% machine-generated captions is preferred
% 70 / 105 (all annotation)
% 22 / 35 (averaging)

\end{comment}


